% concept01.tex
\section{Lines, Angles, and Shapes}
\kr{(선, 각, 도형)}

Geometry begins with the simplest building blocks: points, lines, and angles. A
point marks a position in space. A line extends infinitely in both directions. A
line segment has two endpoints, while a ray starts at one point and goes on
forever in one direction.

An angle forms when two rays meet at a common endpoint called the vertex. We
measure angles in degrees, from 0° to 360°. The most common angles are right
angles (90°), which form perfect corners, and straight angles (180°), which form
straight lines.

\blockquote{A point marks a position, a line extends forever, and an angle forms
when two rays meet at a vertex.}

\subsection{What you'll learn}

You will learn what a point, line, line segment, and ray are. You will also meet
four common types of angles and see basic flat (2D) shapes.

\subsection{Key vocabulary}

\begin{itemize}
  \item \textbf{Point} (\kr{점}) --- a tiny dot that shows a position.
  \item \textbf{Line} (\kr{선}) --- extends forever in both directions.
  \item \textbf{Line segment} (\kr{선분}) --- part of a line with two endpoints.
  \item \textbf{Ray} (\kr{반직선}) --- starts at one point and goes on forever in one direction.
  \item \textbf{Angle} (\kr{각}) --- made when two rays meet at a point (the vertex).
  \item \textbf{Vertex} (\kr{꼭짓점}) --- the meeting point of two sides or rays.
  \item \textbf{Polygon} (\kr{다각형}) --- a shape made of straight sides (triangle, square, etc.).
\end{itemize}

\subsection{Short explanation and examples}

A \emph{point} is a place on paper you can mark with a dot. A \emph{line} keeps
going and never stops; draw arrows on both ends to show this. A \emph{segment} has
two ends; draw a short straight line with dots at both ends. A \emph{ray} has one
end and one arrow.

An \emph{angle} is the space between two rays that share the same endpoint. The
endpoint is the \emph{vertex}. Use a protractor to measure angles.

\subsection{Types of angles (simple list)}

\begin{itemize}
  \item \textbf{Acute} (\kr{예각}): less than \(90^\circ\).
  \item \textbf{Right} (\kr{직각}): exactly \(90^\circ\).
  \item \textbf{Obtuse} (\kr{둔각}): more than \(90^\circ\) and less than \(180^\circ\).
  \item \textbf{Straight} (\kr{평각}): exactly \(180^\circ\) (a straight line).
\end{itemize}

\subsection{Try this (brief activity)}

\begin{enumerate}
  \item Draw a point and label it \(A\).
  \item From \(A\), draw two rays that form a right angle. Label the rays \(AB\) and \(AC\).
  \item Use a protractor to check that the angle \( \angle BAC = 90^\circ\).
\end{enumerate}


% concept02.tex
\section{Triangles and Quadrilaterals}
\kr{(삼각형과 사각형)}

Triangles and quadrilaterals are the most common polygons. Triangles have three
sides and three angles, while quadrilaterals have four sides and four angles.
Each type has special properties that make them useful in geometry and real-world
applications.

Triangles can be classified by their sides (equilateral, isosceles, scalene) or
by their angles (acute, right, obtuse). Quadrilaterals include familiar shapes
like squares, rectangles, and parallelograms, each with their own unique
characteristics.

\blockquote{Triangles have three sides and three angles, while quadrilaterals have
four sides and four angles.}

\subsection{What you'll learn}

You will learn how to classify triangles by their sides and angles, and you will
meet common four-sided shapes (quadrilaterals).

\subsection{Triangles — quick guide}

Triangles have three sides and three angles. Use the lists below to name a triangle:

\paragraph{By sides}
\begin{itemize}
  \item \textbf{Equilateral} (\kr{정삼각형}): three equal sides and three equal angles.
  \item \textbf{Isosceles} (\kr{이등변삼각형}): two equal sides.
  \item \textbf{Scalene} (\kr{부등변삼각형}): no equal sides.
\end{itemize}

\paragraph{By angles}
\begin{itemize}
  \item \textbf{Acute triangle}: all angles \(<90^\circ\).
  \item \textbf{Right triangle}: one angle \(=90^\circ\).
  \item \textbf{Obtuse triangle}: one angle \(>90^\circ\).
\end{itemize}

\subsection{Quadrilaterals — quick guide}

Quadrilaterals have four sides. Common examples:
\begin{itemize}
  \item \textbf{Square} (\kr{정사각형}): four equal sides and four right angles.
  \item \textbf{Rectangle} (\kr{직사각형}): opposite sides equal, four right angles.
  \item \textbf{Rhombus} (\kr{마름모}): four equal sides (angles may not be right).
  \item \textbf{Parallelogram} (\kr{평행사변형}): both pairs of opposite sides are parallel.
  \item \textbf{Trapezoid} (\kr{사다리꼴}): at least one pair of parallel sides.
\end{itemize}

\subsection{Short activity}

\begin{enumerate}
  \item Draw one example of each quadrilateral on grid paper.
  \item For each, list which sides are equal and which angles look equal.
\end{enumerate}


% concept03.tex
\section{Circles and Parts of a Circle}
\kr{(원과 원의 부분)}

Circles are perfectly round shapes where every point on the edge is the same
distance from the center. The radius is the distance from the center to any point
on the circle, while the diameter is twice the radius and passes through the
center.

The circumference is the distance around the circle, and the area is the space
inside. These measurements follow simple formulas using the mathematical constant
$\pi$ (pi), which is approximately 3.14.

\blockquote{A circle is perfectly round with every point on the edge the same
distance from the center.}

\subsection{What you'll learn}

You will learn the names of circle parts and two simple formulas: how to find the
circumference (distance around) and the area (space inside).

\subsection{Key parts}

\begin{itemize}
  \item \textbf{Center} (\kr{중심}): the middle point of a circle.
  \item \textbf{Radius} (\kr{반지름}): distance from center to edge.
  \item \textbf{Diameter} (\kr{지름}): a straight line through the center connecting two points on the circle. \(d = 2r\).
  \item \textbf{Chord} (\kr{현}): any line segment joining two points on the circle.
  \item \textbf{Circumference} (\kr{둘레}): distance around the circle.
  \item \textbf{Area}: space inside the circle.
\end{itemize}

\subsection{Formulas (use simple numbers)}

\[
\text{Circumference} = 2\pi r \quad\text{(or }\pi d\text{)},\qquad
\text{Area} = \pi r^{2}.
\]
(Use \(\pi \approx 3.14\) for easy calculations.)

\subsection{Example}

If a circle has radius \(r=7\) cm, then
\begin{itemize}
  \item Circumference \(= 2\pi r \approx 2\times 3.14\times 7 = 43.96\) cm.
  \item Area \(= \pi r^2 \approx 3.14\times 49 = 153.86\) cm\(^2\).
\end{itemize}


% concept04.tex
\section{Area and Perimeter}
\kr{(넓이와 둘레)}

Area and perimeter are two fundamental measurements of shapes. Area measures the
space inside a shape, while perimeter measures the distance around the outside.
These measurements help us understand how much space shapes take up and how much
material we need to build or cover them.

Different shapes have different formulas for calculating area and perimeter. Some
shapes, like rectangles and triangles, have simple formulas that are easy to
remember and use.

\blockquote{Area measures the space inside a shape, while perimeter measures the
distance around the outside.}

\subsection{What you'll learn}

Find the area (space inside) and the perimeter (distance around) for common flat
shapes.

\subsection{Formulas (memorize the short list)}

\begin{itemize}
  \item \textbf{Rectangle:} \(\text{Area} = \text{length}\times\text{width}\). \\
        \(\text{Perimeter} = 2(\text{length} + \text{width})\).
  \item \textbf{Triangle:} \(\text{Area} = \frac{1}{2}\times\text{base}\times\text{height}\).
  \item \textbf{Circle:} \(\text{Area}=\pi r^{2}\), \(\text{Circumference}=2\pi r\).
\end{itemize}

\subsection{Short practice}

\begin{enumerate}
  \item A rectangle is 8 m long and 3 m wide. Find the area and perimeter.
  \item A triangle has base 6 cm and height 4 cm. Find its area.
\end{enumerate}


% concept05.tex
\section{Surface Area and Volume}
\kr{(겉넓이와 부피)}

Surface area and volume are measurements of three-dimensional objects. Surface
area measures the total area of all the faces or surfaces of an object, while
volume measures how much space the object takes up inside.

These measurements are important in real-world applications like packaging,
construction, and manufacturing. Understanding the difference helps us choose the
right materials and understand how objects relate to space.

\blockquote{Surface area measures the total area of all faces, while volume
measures how much space an object takes up inside.}

\subsection{What you'll learn}

Learn the difference between surface area (sum of all outside faces) and volume
(how much fits inside).

\subsection{Quick formulas}

\begin{itemize}
  \item \textbf{Cube} with side \(s\): \(\text{SA}=6s^{2}\), \(\text{Volume}=s^{3}\).
  \item \textbf{Rectangular prism} \(l\times w\times h\): \(\text{SA}=2(lw+lh+wh)\), \(\text{Volume}=lwh\).
  \item \textbf{Cylinder} radius \(r\), height \(h\): \(\text{SA}=2\pi r^{2}+2\pi r h\), \(\text{Volume}=\pi r^{2}h\).
\end{itemize}

\subsection{Think about it}

Surface area is like the wrapping paper needed to cover a gift box. Volume is how
many small cubes (1×1×1) fit inside the box.


% concept06.tex
\section{Coordinate Plane}
\kr{(좌표평면)}

The coordinate plane is a grid system that helps us locate points and describe
shapes mathematically. It consists of two perpendicular number lines: the x-axis
(horizontal) and the y-axis (vertical). The point where they intersect is called
the origin.

Every point on the plane can be described by an ordered pair (x, y), where x
tells us how far to move horizontally from the origin, and y tells us how far to
move vertically.

\blockquote{The coordinate plane uses two perpendicular axes to locate points with
ordered pairs (x, y).}

\subsection{What you'll learn}

You will place points using ordered pairs \((x,y)\) and draw shapes on the grid.

\subsection{Key ideas}

\begin{itemize}
  \item The horizontal line is the \textbf{x-axis}; the vertical line is the \textbf{y-axis}.
  \item The axes meet at the \textbf{origin} \((0,0)\).
  \item Read an ordered pair as: move \(x\) (left/right), then move \(y\) (up/down).
\end{itemize}

\subsection{Try it}

\begin{enumerate}
  \item Plot these points: \((1,1)\), \((4,1)\), \((4,3)\), \((1,3)\).
  \item Connect the points in order. What shape did you make?
\end{enumerate}


% concept07.tex
\section{Transformations}
\kr{(이동, 회전, 대칭, 확대/축소)}

Transformations are ways to move, rotate, flip, or resize shapes without changing
their basic properties. They help us understand how shapes relate to each other
and are used in computer graphics, engineering, and many real-world applications.

The four basic transformations are translation (sliding), rotation (turning),
reflection (flipping), and dilation (scaling). Each transformation preserves
certain properties of the original shape while changing others.

\blockquote{Transformations move, rotate, flip, or resize shapes while preserving
their basic properties.}

\subsection{What you'll learn}

You will learn four simple moves that change where a shape sits or how big it is:
translation, rotation, reflection, and dilation.

\subsection{Short definitions and examples}

\begin{itemize}
  \item \textbf{\kr{Translation (이동)}}: slide the whole shape; it does not turn or change size. \\
        Example: move every point of a triangle 3 units right and 2 units up.
  \item \textbf{\kr{Rotation (회전)}}: turn the shape around a point (often the origin). \\
        Example: rotate a square \(90^\circ\) clockwise.
  \item \textbf{\kr{Reflection (대칭)}}: flip the shape across a line (like a mirror). \\
        Example: reflect the point \((2,1)\) across the \(y\)-axis to get \((-2,1)\).
  \item \textbf{\kr{Dilation (확대/축소)}}: make the shape bigger or smaller from a center point using a scale factor. \\
        Example: a dilation with factor 2 makes all distances from the center twice as long.
\end{itemize}

\subsection{Classroom activity}

\begin{enumerate}
  \item On graph paper, draw a small shape and then:
    \begin{itemize}
      \item translate it (slide), 
      \item rotate it \(90^\circ\),
      \item reflect it across the \(x\)-axis,
      \item dilate it by factor \(1.5\).
    \end{itemize}
  \item Write one sentence about what changed and what stayed the same after each move.
\end{enumerate}


% ending.tex
\section{Tip for Students}

Geometry is not just about memorizing formulas. It is about understanding how
shapes work and how they relate to the world around us. Learn to visualize
shapes, understand their properties, and see how measurements connect to real
problems. These skills will help you in art, architecture, engineering, and many
other fields.